\begin{proof}[Cauchy–Schwarz inequality in $\mathbb{R}^{n}$]
Let  
\[
\mathbf a=(a_{1},\dots ,a_{n}),\qquad 
\mathbf b=(b_{1},\dots ,b_{n})\in\mathbb R^{n}.
\]
The standard inner product and Euclidean norm are  
\[
\langle\mathbf a,\mathbf b\rangle=\sum_{i=1}^{n}a_i b_i ,\qquad
\|\mathbf a\|=\sqrt{\langle\mathbf a,\mathbf a\rangle}
        =\Bigl(\sum_{i=1}^{n}a_i^{2}\Bigr)^{1/2}.
\]

For any real $t$ consider the non‑negative quadratic form  
\[
f(t)=\|\mathbf a-t\mathbf b\|^{2}
    =\langle\mathbf a-t\mathbf b,\mathbf a-t\mathbf b\rangle
    =\sum_{i=1}^{n}(a_i-t b_i)^{2}\ge 0 .
\]

Expanding the sum gives  
\[
\begin{aligned}
f(t)
   &=\sum_{i=1}^{n}\bigl(a_i^{2}-2t a_i b_i+t^{2}b_i^{2}\bigr)\\
   &=\Bigl(\sum_{i=1}^{n}b_i^{2}\Bigr)t^{2}
     -2\Bigl(\sum_{i=1}^{n}a_i b_i\Bigr)t
     +\sum_{i=1}^{n}a_i^{2}.
\end{aligned}
\]
Set  
\[
A:=\sum_{i=1}^{n}b_i^{2},\qquad
B:=-2\sum_{i=1}^{n}a_i b_i,\qquad
C:=\sum_{i=1}^{n}a_i^{2},
\]
so that $f(t)=At^{2}+Bt+C$.

Since $f(t)\ge0$ for all $t\in\mathbb R$, the quadratic $At^{2}+Bt+C$ cannot have two distinct real roots; consequently its discriminant is non‑positive:
\[
B^{2}-4AC\le 0 .
\]
(If $A=0$ then $\mathbf b=0$, whence $B=0$ and the inequality is trivial.)

Substituting $A,B,C$ yields
\[
\bigl(-2\sum_{i=1}^{n}a_i b_i\bigr)^{2}
   -4\Bigl(\sum_{i=1}^{n}b_i^{2}\Bigr)\Bigl(\sum_{i=1}^{n}a_i^{2}\Bigr)\le0 .
\]
Simplifying,
\[
4\Bigl(\sum_{i=1}^{n}a_i b_i\Bigr)^{2}
   \le 4\Bigl(\sum_{i=1}^{n}a_i^{2}\Bigr)\Bigl(\sum_{i=1}^{n}b_i^{2}\Bigr).
\]
Dividing by $4>0$ we obtain the **Cauchy–Schwarz inequality**
\[
\boxed{\Bigl(\sum_{i=1}^{n}a_i b_i\Bigr)^{2}
       \le\Bigl(\sum_{i=1}^{n}a_i^{2}\Bigr)
          \Bigl(\sum_{i=1}^{n}b_i^{2}\Bigr)} .
\]

\textbf{Equality case.}  
Equality holds iff the discriminant vanishes, i.e. $B^{2}=4AC$.  
In that situation the quadratic $f(t)$ has a double root at
\[
t_{0}= \frac{\sum_{i=1}^{n}a_i b_i}{\sum_{i=1}^{n}b_i^{2}}
      \qquad\bigl(\text{provided }\sum_{i=1}^{n}b_i^{2}\neq0\bigr).
\]
Then $f(t_{0})=0$, so $\mathbf a=t_{0}\mathbf b$. Hence there exists
$\lambda\in\mathbb R$ such that $a_i=\lambda b_i$ for every $i$, i.e.
the two vectors are linearly dependent.  
If $\mathbf b=0$ then both sides of the inequality are zero, and equality
holds trivially; this is consistent with the statement that equality occurs
precisely when the vectors are linearly dependent (the zero vector being
linearly dependent with any vector).

Thus, for any real families $(a_i)_{i=1}^{n}$ and $(b_i)_{i=1}^{n}$,
\[
\left(\sum_{i=1}^{n}a_i b_i\right)^{2}
   \le
\left(\sum_{i=1}^{n}a_i^{2}\right)
\left(\sum_{i=1}^{n}b_i^{2}\right),
\]
with equality exactly when $\mathbf a$ and $\mathbf b$ are linearly
dependent. This is the Cauchy–Schwarz inequality for real sequences.
\end{proof}